% Obsahový zdroj: import_hodiny/31_Faktorial.tex a odpovídající PDF (obě strany).
% P1–P5: zdrojové Příklady 1–5; C1–C5: zdrojové Procvičování 1–5.
% Matematický obsah zachován; výklad, odvození a řešení pedagogicky rozvedeny.
% Upřesnění a opravy: viz KMB_29-41_Kontrola.md; původní import zůstává beze změny.
% Rozpis: 4_rocnik/tematicke_plany/Rozpis_vyuky_MAT_4_rocnik.tex, hodina 31.
% Plán: 4_rocnik/tematicke_plany/SPS.pdf, oddíly 2.1–2.3.
% Sazba: sablony/pripravy/TEMPLATE_Priprava.tex a společný priprava.sty.
% Vlastní doplnění: didaktický výklad, zdůvodnění, metodika; žádné nové číslované úlohy.
% Učebnice/09_Kombinatorika nemá obsahovou kapitolu; místní kombinatorické listy ani Realisticky.cz nejsou dostupné.
\documentclass[11pt,a4paper]{article}
% Spolecny styl pro ucitelske pripravy MAT (Work/MAT)
% Verze 1.1, 2026-10-23
\NeedsTeXFormat{LaTeX2e}
\ProvidesPackage{priprava}[2026/10/23 v1.1 Sablona ucitelskych priprav MAT]

% Zakladni typografie a cestina
\RequirePackage[T1]{fontenc}
\RequirePackage[utf8]{inputenc}
\RequirePackage[czech]{babel}
\RequirePackage{lmodern}
\RequirePackage{microtype}
\RequirePackage{amsmath,amssymb}
\RequirePackage{xcolor}
\RequirePackage{enumitem}
\RequirePackage[most]{tcolorbox}
\RequirePackage{tabularx,booktabs,array}
\RequirePackage{fancyhdr}
\RequirePackage{needspace}
\RequirePackage{geometry}
\RequirePackage[hidelinks]{hyperref}

\geometry{a4paper,top=17mm,bottom=18mm,left=23mm,right=23mm}
\definecolor{notegray}{gray}{0.39}
\setlength{\parindent}{0pt}
\setlength{\parskip}{2.5pt plus 1pt minus 0.5pt}
\setlist[itemize]{topsep=2pt,itemsep=1pt,parsep=0pt,leftmargin=1.45em}
\setlist[enumerate]{topsep=2pt,itemsep=2pt,parsep=0pt,leftmargin=1.65em}
\setcounter{secnumdepth}{0}
\renewcommand{\arraystretch}{1.12}

% Zahlavi: vlevo rocnik | globalni tema, vpravo cislo: nazev hodiny
% Zapati zustava zcela prazdne.
\pagestyle{fancy}
\fancyhf{}
\renewcommand{\headrulewidth}{0.35pt}
\renewcommand{\footrulewidth}{0pt}
\setlength{\headheight}{14pt}
\newcommand{\nastavPripravu}[4]{%
  \fancyhead[L]{\small #1\enspace\textbar\enspace #2}%
  \fancyhead[R]{\small #3: #4}%
}

% Minimalisticke cernobile boxy. Kazda definice zvlast.
\tcbset{commonlesson/.style={enhanced,breakable=false,colback=white,colframe=black,
  colbacktitle=white,coltitle=black,boxrule=0.45pt,
  arc=0.8mm,outer arc=0.8mm,left=3mm,right=3mm,top=1.6mm,bottom=1.6mm,
  toptitle=1mm,bottomtitle=0.7mm,fonttitle=\bfseries,
  attach boxed title to top left={xshift=2.5mm,yshift=-1.15mm},
  boxed title style={colback=white,colframe=white,boxrule=0pt,left=1mm,right=1mm}}}
\newtcolorbox{zadaniBox}[1][]{commonlesson,title={Zadání},#1}
\newtcolorbox{prikladBox}[1][]{commonlesson,title={Příklad},#1}
\newtcolorbox{definiceBox}[1][]{commonlesson,title={Definice},#1}

% Spolecny box pro zadani a jeho reseni, bez mezery mezi castmi.
% Pouziti:
%   \begin{prikladReseniBox}
%     Zadani ulohy.
%     \tcblower
%     \begin{align*} ... \end{align*}
%   \end{prikladReseniBox}
% Spodni cast ma jemny sedy podklad a delici linku; lze v ni sazet
% vice radku matematiky. Obe casti maji jen jeden spolecny ramecek.
\tcbset{solutionlesson/.style={
  commonlesson,
  enhanced jigsaw,
  breakable,
  colbacklower=black!3,
  segmentation style={solid,draw=black!25,line width=0.35pt},
  fontlower=\small\color{notegray},
  before lower={\noindent\textbf{Řešení:}\enspace},
  leftlower=3mm,rightlower=3mm
}}
\newtcolorbox{prikladReseniBox}[1][]{solutionlesson,title={Příklad},#1}
\newtcolorbox{zadaniReseniBox}[1][]{solutionlesson,title={Zadání},#1}

% Pomucky pro autora pripravy: sedive poznamky; \reseni pro starsi pripravy
\newcommand{\poznamka}[1]{\par\begingroup\small\color{notegray}\itshape
  \textnormal{\textbf{Poznámka pro učitele.}}\enspace #1\par\endgroup}
\newcommand{\reseni}[1]{\par\begingroup\small\color{notegray}
  \textbf{Očekávané řešení:}\enspace #1\par\endgroup}
\newcommand{\blok}[1]{\needspace{5\baselineskip}\subsection*{#1}}
\newcommand{\vyrok}[1]{\uv{#1}}

\nastavPripravu{4. ročník}{Kombinatorika}{31}{Faktoriál}
\begin{document}


Připomeň součin počtů možností při postupném obsazování míst. Součiny po sobě jdoucích přirozených čísel se v kombinatorice opakují; potřebujeme pro ně úsporný zápis a pravidla úprav.
\poznamka{Předpoklady: pravidlo součinu, krácení součinů, vytýkání a kvadratická rovnice. Cíl: žák rozumí $0!$, rozvine jen potřebnou část faktoriálu, uvede podmínky a ověří kořeny. Novou látkou není obecná teorie faktoriálových rovnic; P4 a C5 používají pouze kvadratickou rovnici.}

\blok{Výklad a odvození}
\begin{definiceBox}
\textbf{Faktoriál} nezáporného celého čísla $n$, značený $n!$, definujeme předpisem
\[
 n!=1\cdot2\cdots n\quad(n\geq1),\qquad 0!=1.
\]
Symbol $\mathbb N_0$ zde označuje množinu $\{0,1,2,\ldots\}$. Faktoriál záporného ani necelého čísla v této látce nedefinujeme.
\end{definiceBox}

Pro celé $n\geq1$ z definice plyne $n!=n(n-1)!$. Zvláštní hodnota $0!=1$ zajišťuje platnost i pro $n=1$, neboť $1!=1\cdot0!$. Odpovídá také tomu, že prázdný součin má hodnotu jedna; „neudělat žádnou volbu“ je jediná možnost.

\[
\begin{array}{c|rrrrrrr}
 n&0&1&2&3&4&5&6\\\hline
 n!&1&1&2&6&24&120&720
\end{array}
\]
Při krácení rozepiš větší faktoriál až k menšímu. Například $(n+2)!=(n+2)(n+1)n!$ pro $n\in\mathbb N_0$. Faktoriály jsou kladné, a proto je lze v součinu krátit. Obsahuje-li čitatel součet, nejprve vytkneme společný činitel nebo zlomek rozdělíme na součet zlomků.

Faktoriál se \emph{nerozděluje přes součet}: například $(2+3)!=120$, zatímco $2!+3!=8$. Tím vyvracíme obecné pravidlo $(a+b)!=a!+b!$; netvrdíme však, že rovnost nemůže nastat pro některé jednotlivé hodnoty.

U výrazů s proměnnou nejprve stanovíme definiční obor. V této sadě volíme u algebraických úprav $n\in\mathbb N_0$, pokud přítomnost $(n-1)!$ nebo $(n-2)!$ nevynutí užší podmínku. V rovnicích musí splnit podmínky každý argument faktoriálu.


\poznamka{Průběh: 0–8 min definice a $0!$; 8–20 min P1–P3, krácení; 20–30 min P4 a kontrola podmínek, P5 krátce jako kontrast k součinu; 30–42 min C1–C5; 42–45 min kontrola a rekapitulace. Při pomalejším tempu nech C5 k domácímu dokončení, ale řešení s podmínkou prober před hodinou 32.}

\blok{Společně řešené příklady}

% Zdroj P1: import_hodiny/31_Faktorial.tex, Příklad 1.
\begin{prikladReseniBox}[title={Příklad P1}]
Vypočti \(6!\) a \(\frac{8!}{6!}\).
\tcblower
\(6!=1\cdot2\cdot3\cdot4\cdot5\cdot6=720\). Dále \(8!=8\cdot7\cdot6!\), tedy \(\frac{8!}{6!}=8\cdot7=56\).
\end{prikladReseniBox}


% Zdroj P2: import_hodiny/31_Faktorial.tex, Příklad 2.
\begin{prikladReseniBox}[title={Příklad P2}]
Zjednoduš \(\frac{(n+2)!}{n!}\) pro \(n\in\mathbb N_0\).
\tcblower
Pro $n\in\mathbb N_0$ je $n!>0$. Rozvineme pouze dva nejvyšší činitele:
\[
 \frac{(n+2)!}{n!}=\frac{(n+2)(n+1)n!}{n!}=(n+2)(n+1).
\]
Pro krajní $n=0$ vychází na obou stranách $2$, takže podmínka zahrnuje i nulu.
\end{prikladReseniBox}


% Zdroj P3: import_hodiny/31_Faktorial.tex, Příklad 3.
\begin{prikladReseniBox}[title={Příklad P3}]
Vypočti bez kalkulačky \(\frac{9!}{7!\,2!}\).
\tcblower
Čitatel \(9!=9\cdot8\cdot7!\); jmenovatel \(7!\cdot2\). Výsledek \(\frac{9\cdot8}{2}=36\).
\end{prikladReseniBox}


% Zdroj P4: import_hodiny/31_Faktorial.tex, Příklad 4.
\begin{prikladReseniBox}[title={Příklad P4}]
Vyřeš v přirozených číslech \(\frac{(n+1)!}{(n-1)!}=30\).
\tcblower
Musí být $n\in\mathbb N$ a $n\geq1$, aby existoval $(n-1)!$. Pak jsou ekvivalentní úpravy
\begin{align*}
 \frac{(n+1)!}{(n-1)!}=30
 &\iff n(n+1)=30\\
 &\iff n^2+n-30=0\\
 &\iff (n+6)(n-5)=0.
\end{align*}
Kořeny jsou $-6$ a $5$. Záporný kořen je mimo definiční obor. Zkouška pro $n=5$: $6!/4!=6\cdot5=30$. Jediným řešením je $n=5$.
\end{prikladReseniBox}


% Zdroj P5: import_hodiny/31_Faktorial.tex, Příklad 5.
\begin{prikladReseniBox}[title={Příklad P5}]
Pro $n\in\mathbb N_0$ zjednoduš $\dfrac{n!+(n+1)!}{n!}$.
\tcblower
V součtu nejprve vytkneme společný činitel $n!$:
\[
\frac{n!+(n+1)!}{n!}=\frac{n!+(n+1)n!}{n!}
=\frac{n![1+(n+1)]}{n!}=n+2.
\]
Krácení je dovoleno pro každé $n\in\mathbb N_0$, neboť $n!$ není nula.
\end{prikladReseniBox}


\newpage
\blok{Samostatná práce a zásoba k procvičení}

U číselných zlomků nevyčísluj celý čitatel. U C3 pracuj s $n\in\mathbb N_0$; u C5 nejprve zpřesni obor podle jmenovatele.

% Zdroj C1: import_hodiny/31_Faktorial.tex, Procvičování, položka 1.
\begin{zadaniReseniBox}[title={Úloha C1}]
Vypočti \(5!\).
\tcblower
Postupným násobením $5!=1\cdot2\cdot3\cdot4\cdot5=120$.
\end{zadaniReseniBox}

% Zdroj C2: import_hodiny/31_Faktorial.tex, Procvičování, položka 2.
\begin{zadaniReseniBox}[title={Úloha C2}]
Zjednoduš \(\frac{11!}{9!}\).
\tcblower
$11!=11\cdot10\cdot9!$, tedy $\dfrac{11!}{9!}=11\cdot10=110$.
\end{zadaniReseniBox}

% Zdroj C3: import_hodiny/31_Faktorial.tex, Procvičování, položka 3.
\begin{zadaniReseniBox}[title={Úloha C3}]
Zjednoduš $\dfrac{(n+3)!}{(n+1)!}$ pro $n\in\mathbb N_0$.
\tcblower
Pro zadané $n$ oba faktoriály existují a jmenovatel je kladný. Rozvineme pouze potřebnou část:
\[
\frac{(n+3)!}{(n+1)!}
=\frac{(n+3)(n+2)(n+1)!}{(n+1)!}=(n+3)(n+2).
\]
\end{zadaniReseniBox}

% Zdroj C4: import_hodiny/31_Faktorial.tex, Procvičování, položka 4.
\begin{zadaniReseniBox}[title={Úloha C4}]
Vypočti \(\frac{7!}{5!\,2!}\).
\tcblower
$\dfrac{7!}{5!\,2!}=\dfrac{7\cdot6\cdot5!}{5!\cdot2}=21$.
\end{zadaniReseniBox}

% Zdroj C5: import_hodiny/31_Faktorial.tex, Procvičování, položka 5.
\begin{zadaniReseniBox}[title={Úloha C5}]
Urči \(n\in\mathbb N\), pro které \(\frac{n!}{(n-2)!}=42\).
\tcblower
Podmínka je celé $n\geq2$. Po zkrácení dostaneme $n(n-1)=42$, tedy $(n-7)(n+6)=0$. Kořen $-6$ nevyhovuje podmínce, $n=7$ ano. Zkouška: $7!/5!=7\cdot6=42$. Ekvivalentní kvadratická rovnice zajišťuje, že jiné řešení není.
\end{zadaniReseniBox}


\poznamka{Typické chyby: $0!=0$; krácení jednotlivých sčítanců místo společného činitele; rozvoj faktoriálu až k jedničce; přijetí záporného kořene. Nech žáka u P4 slovně vysvětlit, proč $n=-6$ nepatří do zadání. Příště součin $n(n-1)\cdots1$ odvodíme jako počet uspořádání všech prvků.}
\end{document}
